% Generated draft: requires mathematical and source review.
\documentclass[11pt]{article}
\usepackage[margin=25mm]{geometry}
\usepackage{fontspec}
\setmainfont{Latin Modern Roman}
\usepackage{amsmath,amssymb,mathtools,unicode-math}
\setmathfont{Latin Modern Math}
\usepackage{enumitem,needspace,xurl,hyperref}
\setlength{\parindent}{0pt}
\setlength{\parskip}{5pt}
\setlength{\emergencystretch}{3em}
\newcommand{\N}{\mathbb{N}}
\newcommand{\Z}{\mathbb{Z}}
\newcommand{\Q}{\mathbb{Q}}
\newcommand{\R}{\mathbb{R}}
\newcommand{\C}{\mathbb{C}}
\newcommand{\F}{\mathbb{F}}
\newcommand{\A}{\mathbb{A}}
\newcommand{\PP}{\mathbb{P}}
\newcommand{\E}{\mathbb{E}}
\newcommand{\eps}{\varepsilon}
\newcommand{\dd}{\,\mathrm{d}}
\begin{document}
\section*{UnsolvedMath ID: 30000201}
\textbf{Source problem code:} catalogue entry
\subsection{Definitions and mathematical statement}
The record proposes two implications, with no restriction to finite polygons:
\[
 \forall\Gamma,\quad
 \Gamma\text{ is a }2\text{-Moufang generalized }n\text{-gon}
 \ \Longrightarrow\ \Gamma\text{ is Moufang},
\]
\[
 \forall\Gamma,\quad
 \Gamma\text{ is a half-Moufang polygon}
 \ \Longrightarrow\ \Gamma\text{ is Moufang}.
\]
Each quantifier ranges over the polygons and admissible integers \(n\) intended in the report. The assertions are separate; their combined hypothesis is a disjunction, not a requirement that a polygon satisfy both weaker conditions.

For orientation, the usual incidence-graph definition of a generalized \(n\)-gon is a connected bipartite point-line graph with diameter \(n\) and girth \(2n\). Thickness means that each vertex has at least three neighbors. Apartments are cycles of length \(2n\), and roots are paths of length \(n\) in apartments. In the usual root-group convention, for a root \(\alpha=(v_0,\ldots,v_n)\), its root group consists of the type-preserving automorphisms fixing every vertex adjacent to each of \(v_1,\ldots,v_{n-1}\). A polygon is Moufang when each such root group acts transitively on the apartments containing its root. These standard conventions explain the conclusion; the supplied record does not specify thickness, allowable small values of \(n\), or whether an ambient automorphism group rather than the full automorphism group is part of its setup.

Crucially, the definitions of \(2\)-Moufang and half-Moufang are not supplied. They must be recovered with their exact local transitivity, root, type and group conventions before either implication is a completely self-contained theorem. No definition based on only two roots, a rank-two building, or one arbitrarily chosen half of the roots is guessed from the names. No classification of finite or infinite generalized polygons is added as an unstated hypothesis.

The literature assessment checked 2026-08-21 labels the item partially solved: it reports the \(2\)-Moufang implication proved for generalized \(n\)-gons with \(n\leq6\), while the unrestricted infinite case and the half-Moufang implication remain open. This is retained as a dated report of progress, with the theorem's precise hypotheses absent. It is not promoted to an independently checked theorem or an assertion that every larger finite case is unresolved. The source-verification field reports checking and normalizing the official statement but does not provide the missing definitions.
\subsection{Short English statement}
The conjecture asks separately whether the \(2\)-Moufang condition and the half-Moufang condition each force a polygon to be Moufang. The source reports partial progress for the first implication through \(n\leq6\), with infinite and half-Moufang cases still unresolved in its dated assessment. Exact definitions of the weaker conditions and the polygon conventions are missing.
\subsection{Sources}
\begin{itemize}
\item Local UnsolvedMath registry, permanent ID 30000201, source code catalogue entry, ``Moufangness of Two-Moufang and Half-Moufang Polygons,'' including original conjecture, cleaned question and verification metadata.
\item \emph{Groups and Geometries (2005)}, Oberwolfach Reports, pp.\ 647--721, DOI 10.4171/owr/2005/12: \url{https://doi.org/10.4171/owr/2005/12}, supplied as the ``original conjecture.''
\item Supplied reference ``Tent theorem for n at most 6'': \url{https://doi.org/10.1016/j.ejc.2004.06.017}.
\item Supplied reference ``survey recording remaining cases'': \url{https://cage.ugent.be/~hvm/artikels/107.pdf}. The URL is cited as supplied; no PDF has been opened.
\item Literature triage checked 2026-08-21; local record updated 2026-09-16. No external source has been visited or independently checked here.
\end{itemize}
\end{document}
